\section{Qualifications of the PIs}
\label{sec:pis}

PIs Kruegel and Vigna have a twenty-year history of collaboration.
They have developed novel approaches to vulnerability analysis and
malware detection and have created tools to combat cybercrime and
abuse of social networks. The results of this collaboration include
widely-adopted open-source tools (such as \texttt{angr} and earlier
systems like Anubis), datasets shared with the community, and
hundreds of peer-reviewed papers across a broad range of topics in
security and privacy. Both PIs have substantial experience managing
large research grants from NSF, DARPA, and other agencies, and
co-lead a large lab with dozens of graduate and undergraduate students.

The PIs' experience is directly relevant to each of the three thrusts
of this proposal. Their work on modeling opaque systems from observed
interactions (HALucinator, Conware, and earlier efforts)---originally
motivated by firmware analysis but methodologically general---provides
the foundation for the opaque-component replication portion of
Thrust~2. Their work on AI planning for multi-step attacks
(ChainReactor) provides the foundation for Thrust~3, and their
participation in the DARPA Cyber Grand Challenge and the DARPA AIxCC
competition has given them direct experience building autonomous
systems that identify and exploit vulnerabilities in complex software,
culminating most recently in the Artiphishell system.

Beyond technical qualifications, the PIs bring a long record of
producing artifacts that are used outside of academia: the open-source
tools noted above; the Lastline company that commercialized earlier
research; and the iCTF competition that has trained cohorts of
security professionals since 2003. The combination of deep technical
expertise, sustained open-source engagement, and established industry
relationships positions the team well to execute the proposed work
and ensure that its outputs reach both the research community and the
operators of the critical infrastructure that \sysname is designed to
help defend.
